\section{Заключение}

В рамках курсового проекта была разработана веб-версия бизнес-игры «Мультикомандный Featureban», предназначенная для моделирования совместной работы нескольких команд в условиях ограничений по WIP, взаимных зависимостей и необходимости регулярной синхронизации состояния. В ходе выполнения работы была подтверждена возможность переноса логики настольной обучающей симуляции в онлайн-формат без потери ключевых образовательных свойств игры.

В процессе проектирования были проанализированы предметная область, сетевые протоколы, архитектурные подходы к построению клиентской и серверной частей, а также методологии организации разработки. На основе проведенного анализа была выбрана архитектура на базе SPA, реализована серверная часть на языке Go, обеспечено взаимодействие клиентов и сервера в реальном времени посредством WebSocket, а также предусмотрен резервный механизм HTTP Polling для повышения устойчивости системы при нарушении основного канала связи.

Отдельное внимание в проекте было уделено логике игровой симуляции: механике досок команд, перемещению задач, ограничению WIP, ретроспективным изменениям параметров процесса и автоматическому переходу между фазами игрового цикла. Реализованные решения позволили формализовать правила игры, сделать их воспроизводимыми в цифровой среде и обеспечить одновременную работу нескольких пользователей в рамках одной сессии.

Практическая значимость разработанного приложения заключается в возможности его использования для дистанционных корпоративных тренингов, учебных занятий и демонстрации принципов Kanban-подхода в мультикомандной среде. Полученный результат может быть использован как основа для дальнейшего развития продукта, включая расширение аналитических возможностей, улучшение пользовательского интерфейса, развитие механизмов отказоустойчивости и добавление новых сценариев взаимодействия между командами.
